\chapter{Schema Reference}
\label{ch:schema-reference}

This chapter is the authoritative reference for every schema construct in the \texttt{gert.ops}
Domain Kit. It covers the Kit declaration, all five \texttt{gert.ops} step types, runbook
metadata fields, and the three evidence types. For the core gert v2 schema (step types
\texttt{cli}, \texttt{manual}, \texttt{tool}, \texttt{branch}, \texttt{iterate}), refer to the
\emph{gert v2 Design Document}, §3.

% ─────────────────────────────────────────────────────────────────────────────
\section{Kit Declaration and \texttt{apiVersion}}
\label{sec:schema-apiversion}
% ─────────────────────────────────────────────────────────────────────────────

A \texttt{gert.ops} runbook declares both the core schema version and the Kit in the
\texttt{apiVersion} and \texttt{kit} fields at the top of the file.

\begin{minted}{yaml}
apiVersion: runbook/v2
kit: ops/v1

meta:
  name: my-ops-runbook
  kind: change
  description: "Short description of what this runbook does."
\end{minted}

\begin{description}
  \item[\texttt{apiVersion: runbook/v2}] Required. Selects the gert v2 core schema. The
    \texttt{gert.ops} Kit is not compatible with \texttt{runbook/v1} runbooks.
  \item[\texttt{kit: ops/v1}] Required for \texttt{gert.ops} runbooks. Activates the Kit
    compiler, making \texttt{ops.*} step types and metadata fields available. Omitting this
    field produces a validation error when any \texttt{ops.*} type is used.
\end{description}

% ─────────────────────────────────────────────────────────────────────────────
\section{Runbook Metadata Fields}
\label{sec:schema-metadata}
% ─────────────────────────────────────────────────────────────────────────────

The \texttt{meta} block in a \texttt{gert.ops} runbook extends the core metadata with
Kit-specific fields. Core fields (\texttt{name}, \texttt{kind}, \texttt{description},
\texttt{inputs}, \texttt{governance}) continue to work as described in the core documentation.

\subsection{\texttt{meta.roles}}

Declares the named participants for this runbook execution. Each field accepts a string (a
specific username or on-call alias) or a list of strings.

\begin{minted}{yaml}
meta:
  roles:
    dri: "@alice"
    approver:
      - "@bob"
      - "@change-board"
    change-manager: "@change-board"
    responder:
      - "@oncall-sre"
    incident-commander: "@ic-oncall"
    observer:
      - "@audit-team"
\end{minted}

\begin{description}
  \item[\texttt{dri}] The Directly Responsible Individual for this run. Exactly one value.
    If omitted, the person starting the run becomes the implicit DRI.
  \item[\texttt{approver}] One or more identities that may approve \texttt{ops.approval}
    gates. Gate-level \texttt{approvers} fields override this list.
  \item[\texttt{change-manager}] The identity responsible for change authorization.
    Required when the runbook contains an \texttt{ops.change-request} step.
  \item[\texttt{responder}] Identities that may execute \texttt{ops.manual} steps during
    an incident. Required when the runbook contains an \texttt{ops.incident} wrapper.
  \item[\texttt{incident-commander}] The identity that assumes DRI during incident
    declaration. Defaults to \texttt{dri} if unset.
  \item[\texttt{observer}] Read-only participants who receive run notifications and can view
    the evidence bundle but cannot take action.
\end{description}

\subsection{\texttt{meta.change-id}}

A reference to an external change ticket (e.g., ServiceNow, Jira) that authorized this
runbook execution. Used for audit trail cross-referencing.

\begin{minted}{yaml}
meta:
  change-id: "CHG0042817"
\end{minted}

If the runbook contains an \texttt{ops.change-request} step, \texttt{meta.change-id} is
populated automatically from the change request record after pre-execution approval is granted.

\subsection{\texttt{meta.sla}}

Defines SLA timeouts that apply at the runbook level. These are defaults; individual steps
may override them.

\begin{minted}{yaml}
meta:
  sla:
    approval-timeout: 30m
    execution-timeout: 4h
    escalation-target: "@sre-lead"
\end{minted}

\begin{description}
  \item[\texttt{approval-timeout}] Maximum time to wait for a gate approval before the run
    is automatically escalated. Duration string (e.g., \texttt{15m}, \texttt{1h}).
  \item[\texttt{execution-timeout}] Maximum total wall-clock time for the entire runbook run.
    If this duration elapses, the run is suspended and the DRI is notified.
  \item[\texttt{escalation-target}] Identity notified when a timeout fires.
\end{description}

% ─────────────────────────────────────────────────────────────────────────────
\section{Step Types}
\label{sec:schema-step-types}
% ─────────────────────────────────────────────────────────────────────────────

The \texttt{gert.ops} Kit defines five step types. All five are lowered to core primitives at
compile time. The runtime never executes \texttt{ops.*} nodes directly.

\subsection{\texttt{ops.cli}}

An audited command-line step. Extends the core \texttt{cli} type with an explicit command
allowlist, argument-level redaction, and an automatic evidence capture option.

\begin{minted}{yaml}
- step:
    id: deploy_image
    type: ops.cli
    title: "Deploy container image"
    command: kubectl
    args:
      - rollout
      - restart
      - "deployment/{{ .service_name }}"
      - "-n"
      - "{{ .namespace }}"
    allowed_commands: [kubectl]
    redact:
      - pattern: "--token=[^ ]+"
        replace: "--token=[REDACTED]"
    capture:
      rollout_output: stdout
    evidence:
      auto: true
      label: "kubectl rollout output"
\end{minted}

\begin{description}
  \item[\texttt{command}] The executable to run. Subject to the governance allowlist.
  \item[\texttt{args}] Argument list. Template expressions are evaluated before execution.
  \item[\texttt{allowed\_commands}] Step-level allowlist override. Merges with the runbook-level
    \texttt{meta.governance.allowed\_commands}.
  \item[\texttt{redact}] List of RE2 pattern/replace pairs applied to stdout and stderr before
    capture and evidence writing.
  \item[\texttt{capture}] Maps run variables to captured output streams.
  \item[\texttt{evidence.auto}] When \texttt{true}, the step output is automatically captured
    as an \texttt{ops.evidence.command-output} record.
  \item[\texttt{evidence.label}] Human-readable label for the auto-captured evidence record.
\end{description}

\subsection{\texttt{ops.manual}}

A manual procedure step requiring human action and evidence submission. Extends the core
\texttt{manual} type with a role requirement and structured evidence specification.

\begin{minted}{yaml}
- step:
    id: verify_deployment
    type: ops.manual
    title: "Verify deployment is healthy"
    requires_role: dri
    instructions: |
      1. Open the Grafana dashboard for {{ .service_name }}.
      2. Confirm error rate < 0.1% over the last 5 minutes.
      3. Confirm p99 latency < 200ms.
      4. Take a screenshot of the dashboard.
    evidence:
      - type: ops.evidence.screenshot
        label: "Grafana dashboard post-deploy"
        required: true
      - type: ops.evidence.attestation
        label: "DRI sign-off"
        required: true
        prompt: "I confirm the service is healthy and the deployment is successful."
\end{minted}

\begin{description}
  \item[\texttt{requires\_role}] The role that must execute this step. One of: \texttt{dri},
    \texttt{approver}, \texttt{change-manager}, \texttt{responder}, \texttt{incident-commander}.
  \item[\texttt{instructions}] Free-form instructions shown to the executor. Template expressions
    are supported.
  \item[\texttt{evidence}] List of evidence specifications. Each entry has \texttt{type},
    \texttt{label}, \texttt{required}, and type-specific fields (see
    §\ref{sec:schema-evidence}).
\end{description}

\subsection{\texttt{ops.approval}}

A blocking gate that requires explicit approval from one or more designated approvers before
execution may continue. Lowers to a \texttt{manual} step with an approval-specific evidence
record and an SLA timeout.

\begin{minted}{yaml}
- step:
    id: change_board_approval
    type: ops.approval
    title: "Change board sign-off"
    description: |
      Review the deployment plan and approve to proceed.
      Runbook: {{ .runbook_name }}
      Target environment: {{ .env }}
      Scheduled window: {{ .change_window }}
    approvers:
      - "@change-board"
      - "@sre-lead"
    requires_any: true
    timeout: 24h
    on_timeout: escalate
    escalation_target: "@vp-engineering"
    evidence:
      - type: ops.evidence.attestation
        label: "Change board approval"
        required: true
        prompt: "I approve this change for the scheduled window."
\end{minted}

\begin{description}
  \item[\texttt{description}] Context shown to approvers in the approval request. Templates
    are evaluated when the gate is reached.
  \item[\texttt{approvers}] List of identities that may sign this gate. Overrides
    \texttt{meta.roles.approver} for this specific gate.
  \item[\texttt{requires\_any}] When \texttt{true} (default), any single approver may sign.
    When \texttt{false}, \emph{all} listed approvers must sign.
  \item[\texttt{timeout}] Duration to wait for approval. Defaults to
    \texttt{meta.sla.approval-timeout} if not set.
  \item[\texttt{on\_timeout}] Action when timeout fires. \texttt{escalate} notifies the
    \texttt{escalation\_target} and suspends the run; \texttt{abort} cancels the run.
  \item[\texttt{escalation\_target}] Identity notified on timeout. Overrides
    \texttt{meta.sla.escalation-target}.
\end{description}

\subsection{\texttt{ops.change-request}}

A wrapper step that groups a deployment or rollback sequence under a formal change request
lifecycle. Provides pre-execution approval, rollback injection on failure, and automatic
change record closure.

\begin{minted}{yaml}
- step:
    id: deploy_change
    type: ops.change-request
    title: "Deploy v2.4.1 to production"
    change-id: "{{ .change_ticket }}"
    risk: low
    rollback:
      runbook: rollback-service.runbook.yaml
      inputs:
        service_name: "{{ .service_name }}"
        target_version: "{{ .previous_version }}"
    steps:
      - step:
          id: pull_image
          type: ops.cli
          title: "Pull container image"
          command: docker
          args: [pull, "{{ .image_ref }}"]
          evidence:
            auto: true
      - step:
          id: apply_manifests
          type: ops.cli
          title: "Apply Kubernetes manifests"
          command: kubectl
          args: [apply, -f, "{{ .manifest_path }}"]
          evidence:
            auto: true
\end{minted}

\begin{description}
  \item[\texttt{change-id}] Reference to the external change ticket. Populates
    \texttt{meta.change-id} if not already set.
  \item[\texttt{risk}] Risk classification: \texttt{low}, \texttt{medium}, \texttt{high},
    \texttt{emergency}. Used by the change-manager gate to determine required approvals.
  \item[\texttt{rollback}] Rollback specification. If any step within \texttt{steps} fails,
    gert automatically initiates the specified rollback runbook.
  \item[\texttt{rollback.runbook}] Path to the rollback runbook (relative to the current
    runbook file).
  \item[\texttt{rollback.inputs}] Input bindings passed to the rollback runbook. Template
    expressions reference the current run's variables.
  \item[\texttt{steps}] The change execution steps, executed in order.
\end{description}

\subsection{\texttt{ops.incident}}

A wrapper step that declares an operational incident and activates incident-response mode.
All steps within the wrapper are executed under incident-commander authority, with mandatory
evidence capture and automatic SLA escalation.

\begin{minted}{yaml}
- step:
    id: declare_incident
    type: ops.incident
    title: "Service degradation — payment processing"
    severity: sev2
    incident-id: "INC-{{ .incident_ticket }}"
    commander: "@ic-oncall"
    sla:
      mitigation-timeout: 30m
      resolution-timeout: 4h
      escalation-target: "@vp-engineering"
    steps:
      - step:
          id: initial_triage
          type: ops.manual
          title: "Initial triage"
          requires_role: responder
          instructions: |
            Check error rates, latency, and recent deployments.
          evidence:
            - type: ops.evidence.attestation
              label: "Triage findings"
              required: true
              prompt: "Describe the symptoms and initial findings."
\end{minted}

\begin{description}
  \item[\texttt{severity}] Incident severity: \texttt{sev1}, \texttt{sev2}, \texttt{sev3},
    \texttt{sev4}. Severity \texttt{sev1} requires immediate IC takeover.
  \item[\texttt{incident-id}] External incident ticket reference (e.g., PagerDuty, Opsgenie).
  \item[\texttt{commander}] Identity that assumes DRI for the duration of the incident wrapper.
    Defaults to \texttt{meta.roles.incident-commander}.
  \item[\texttt{sla.mitigation-timeout}] Time allowed to achieve initial mitigation.
  \item[\texttt{sla.resolution-timeout}] Total time allowed for full resolution.
  \item[\texttt{steps}] Incident response steps executed under incident-commander authority.
\end{description}

% ─────────────────────────────────────────────────────────────────────────────
\section{Evidence Types}
\label{sec:schema-evidence}
% ─────────────────────────────────────────────────────────────────────────────

\texttt{gert.ops} defines three evidence types. Evidence records are written to the run trace
and included in the evidence bundle. All evidence records include a timestamp, the submitting
user's identity, and their role at time of submission.

\subsection{\texttt{ops.evidence.screenshot}}

Captures a screenshot or image file submitted by the executor.

\begin{minted}{yaml}
evidence:
  - type: ops.evidence.screenshot
    label: "Grafana dashboard — post-deploy"
    required: true
    description: "Screenshot showing error rate and latency for 5 minutes post-deploy."
\end{minted}

\begin{description}
  \item[\texttt{label}] Short identifier for this evidence record in the bundle.
  \item[\texttt{required}] When \texttt{true}, the step cannot complete without this evidence.
  \item[\texttt{description}] Instructions shown to the submitter explaining what to capture.
\end{description}

At submission time, gert hashes the uploaded file (SHA-256), stores the binary under
\texttt{.runbook/runs/<run-id>/attachments/<sha256>.bin}, and writes an
\texttt{evidence/submitted} trace event with the hash and label.

\subsection{\texttt{ops.evidence.command-output}}

Captures the stdout/stderr of a CLI command as an evidence record. Produced automatically
by \texttt{ops.cli} steps with \texttt{evidence.auto: true}, or explicitly declared.

\begin{minted}{yaml}
evidence:
  - type: ops.evidence.command-output
    label: "kubectl rollout status"
    required: true
    capture_from: rollout_output
\end{minted}

\begin{description}
  \item[\texttt{capture\_from}] Run variable containing the output to capture. Must be
    populated by a preceding \texttt{capture} block.
\end{description}

Command output is passed through the runbook's \texttt{redact} rules before storage.

\subsection{\texttt{ops.evidence.attestation}}

A structured sign-off where the executor explicitly affirms a statement. Used for DRI
sign-offs, approvals, and change-manager authorizations.

\begin{minted}{yaml}
evidence:
  - type: ops.evidence.attestation
    label: "DRI deployment sign-off"
    required: true
    prompt: "I confirm this deployment is complete and the service is healthy."
\end{minted}

\begin{description}
  \item[\texttt{prompt}] The exact statement the executor affirms. Displayed verbatim to
    the submitter before they confirm.
\end{description}

The attestation record stores the submitter's identity, their role, the prompt text, and the
timestamp. It cannot be submitted by a user who does not hold the step's \texttt{requires\_role}.
\newline

\noindent\textbf{Evidence Record Written to Trace:}
\begin{minted}{yaml}
# Trace event: evidence/submitted
event_id: "4a7b2c91-..."
kind: evidence/submitted
timestamp: "2026-05-14T03:22:11Z"
payload:
  step_id: "verify_deployment"
  evidence_label: "DRI deployment sign-off"
  evidence_type: "ops.evidence.attestation"
  submitted_by: "@alice"
  role: "dri"
  prompt: "I confirm this deployment is complete and the service is healthy."
  sha256: "e3b0c44298fc1c149afbf4c8996fb92427ae41e4649b934ca495991b7852b855"
\end{minted}
